\chapter{PHÂN TÍCH YÊU CẦU HỆ THỐNG VÀ KIỂM SOÁT PHÂN QUYỀN}

\section{Phân tích các vai trò người dùng và phạm vi nghiệp vụ}

Trong giai đoạn phân tích KT1, khi tham vấn công cụ AI về các tác nhân của hệ thống, AI từng gợi ý 5 vai trò bao gồm cả Khách hàng đăng nhập. Tuy nhiên, đối chiếu với thực tế vận hành tại xưởng và phạm vi đề tài 80277, nhóm quyết định chỉ giữ \textbf{4 vai trò nhân sự đăng nhập nội bộ} trong bảng \texttt{users} (\texttt{MANAGER}, \texttt{RECEPTIONIST}, \texttt{TECHNICIAN}, \texttt{CASHIER}), còn Khách hàng được quản lý như một thực thể dữ liệu tại bảng \texttt{customers}. Chi tiết mục tiêu và giới hạn quyền của từng vai trò được trình bày tại Bảng~\ref{tab:g-actors} và sơ đồ Use Case tại Hình~\ref{fig:g-usecase}.

\begin{table}[H]
\centering
\small
\caption{Bốn vai trò người dùng, tài khoản demo và giới hạn quyền hạn trong hệ thống}
\label{tab:g-actors}
\begin{tabularx}{\textwidth}{|L{3.2cm}|Y|Y|}
\hline
\thead{Vai trò \& Tài khoản} & \thead{Nhiệm vụ nghiệp vụ chính} & \thead{Giới hạn quyền hạn (\texttt{security\_matrix.py})}\\
\hline
Quản lý xưởng (\texttt{MANAGER}\\TK: \texttt{quan\_ly}) & Quản trị người dùng (\texttt{user.manage}), sửa danh mục dịch vụ/phụ tùng (\texttt{catalog.edit}), duyệt báo giá (\texttt{quote.approve}), lập hóa đơn (\texttt{invoice.manage}), duyệt kết quả AI (\texttt{ai.review}) và xem báo cáo. & Toàn quyền điều hành; là vai trò duy nhất có quyền phê duyệt hoặc từ chối kết quả AI (\texttt{APPROVED}/\texttt{REJECTED}) và duyệt báo giá chính thức. \\
\hline
Lễ tân dịch vụ (\texttt{RECEPTIONIST}\\TK: \texttt{le\_tan}) & Quản lý khách hàng, phương tiện, lịch hẹn; tiếp nhận xe và lập biên bản (\texttt{intake\_records}); phân công KTV; lập báo giá nháp; nhập/điều chỉnh kho; tạo bảo hành; chạy AI01, AI02, AI03. & Được gọi cả 3 chức năng AI (\texttt{ai.history}, \texttt{ai.explain}, \texttt{ai.quote}) nhưng \textbf{bị chặn} quyền duyệt kết quả AI (\texttt{ai.review}) và duyệt báo giá (\texttt{quote.approve}). \\
\hline
Kỹ thuật viên (\texttt{TECHNICIAN}\\TK: \texttt{ky\_thuat\_vien}) & Xem phiếu sửa chữa, cập nhật trạng thái kiểm tra/sửa chữa (\texttt{order.progress}), thêm/gỡ hạng mục dịch vụ và phụ tùng (\texttt{order.item}), ghi nhật ký kỹ thuật (\texttt{repair\_logs}) và gọi AI03 (\texttt{ai.quote}). & \textbf{Cô lập theo phiếu được giao (\texttt{duoc\_gan\_cong})}: Chỉ được sửa trạng thái, thêm hạng mục, ghi nhật ký hoặc chạy AI03 trên phiếu có \texttt{technician\_id == user.id}. \\
\hline
Thu ngân (\texttt{CASHIER}\\TK: \texttt{thu\_ngan}) & Tra cứu phiếu sửa chữa và hóa đơn, ghi nhận thanh toán từng phần hoặc toàn bộ (\texttt{payment.record}) và hủy khoản thu nhầm. & Chỉ có quyền ghi ở phân hệ Thanh toán; bị chặn (HTTP 403) khi cố tạo/sửa phiếu sửa chữa, chỉnh kho hoặc chạy AI. \\
\hline
\end{tabularx}
\end{table}

\begin{figure}[H]
\centering
\vspace{2pt}
\resizebox{0.98\textwidth}{!}{%
\begin{tikzpicture}[
  font=\sffamily,
  actor/.style={rectangle, rounded corners=3pt, draw=black, line width=0.7pt, fill=black!10, text width=2.8cm, minimum height=0.95cm, align=center, font=\sffamily\scriptsize\bfseries, inner sep=3pt},
  uc/.style={rectangle, rounded corners=8pt, draw=black, line width=0.65pt, fill=white, text width=4.0cm, minimum height=0.68cm, align=center, font=\sffamily\scriptsize, inner sep=3pt},
  aiuc/.style={rectangle, rounded corners=8pt, draw=black, line width=0.75pt, fill=black!8, text width=4.0cm, minimum height=0.68cm, align=center, font=\sffamily\scriptsize\bfseries, inner sep=3pt},
  arr/.style={-{Latex[length=2mm,width=1.4mm]}, line width=0.55pt, draw=black}
]
  \draw[rounded corners=5pt, line width=0.8pt, draw=black] (-2.9,-3.7) rectangle (7.8,3.8);
  \node[font=\sffamily\footnotesize\bfseries, fill=white, inner sep=3pt] at (2.45, 3.8) {RANH GIỚI NGHIỆP VỤ GARAGETOOLS AI (ĐỀ TÀI 80277)};

  \node[uc] (uc1) at (-0.5, 2.7) {UC01: Quản lý Khách hàng \& Xe (VIN)};
  \node[uc] (uc2) at (-0.5, 1.55) {UC02: Đặt lịch hẹn \& Tiếp nhận xe};
  \node[uc] (uc3) at (-0.5, 0.4) {UC03: Lập \& Điều phối Phiếu sửa chữa (RO)};
  \node[uc] (uc4) at (-0.5, -0.75) {UC04: Quản lý Kho \& Đặt giữ phụ tùng};
  \node[uc] (uc5) at (-0.5, -1.9) {UC05: Xuất Hóa đơn, Thu tiền \& Công nợ};
  \node[uc] (uc6) at (-0.5, -3.05) {UC06: Thống kê KPI \& Sao lưu CSDL};

  \node[aiuc] (ai1) at (4.9, 2.5) {AI01: Chẩn đoán triệu chứng kỹ thuật};
  \node[aiuc] (ai2) at (4.9, 1.15) {AI02: Đề xuất hạng mục \& phụ tùng};
  \node[aiuc] (ai3) at (4.9, -0.2) {AI03: Tóm tắt tiến độ sửa chữa RO};
  \node[aiuc] (ai4) at (4.9, -1.55) {AI04: Cảnh báo bảo dưỡng định kỳ};
  \node[uc] (hitl) at (4.9, -2.9) {HITL: Duyệt / Từ chối đề xuất AI};

  \node[actor] (sa) at (-5.2, 2.0) {Cố vấn dịch vụ\\(Service Advisor)};
  \node[actor] (ktv) at (-5.2, -0.2) {Kỹ thuật viên\\(Technician)};
  \node[actor] (kt) at (-5.2, -2.4) {Kế toán / Quản lý\\(Accountant / Admin)};

  \draw[arr] (sa.east) -- (uc1.west);
  \draw[arr] (sa.east) -- (uc2.west);
  \draw[arr] (sa.east) -- (uc3.west);
  \draw[arr] (ktv.east) -- (uc3.west);
  \draw[arr] (ktv.east) -- (uc4.west);
  \draw[arr] (kt.east) -- (uc5.west);
  \draw[arr] (kt.east) -- (uc6.west);

  \draw[arr, dashed] (uc3.east) -- (ai1.west);
  \draw[arr, dashed] (uc3.east) -- (ai2.west);
  \draw[arr, dashed] (uc3.east) -- (ai3.west);
  \draw[arr, dashed] (uc1.east) -- (ai4.west);
  \draw[arr] (hitl.east) -- ++(0.35, 0) |- (ai2.east);
\end{tikzpicture}%
}

\vspace{2pt}
\caption{Sơ đồ ca sử dụng tổng thể (Use Case Diagram) của hệ thống GarageTools AI}
\label{fig:g-usecase}
\end{figure}

\section{Danh mục yêu cầu chức năng (FR)}

Hệ thống hiện thực hóa đầy đủ 14 nhóm yêu cầu chức năng (FR01--FR14) bao phủ toàn bộ quy trình vận hành garage và 3 chức năng trợ lý AI, được tổng hợp cùng kết quả kiểm chứng trong Bảng~\ref{tab:g-functional}.

\begin{small}
\renewcommand{\arraystretch}{1.08}
\begin{longtable}{|C{1.0cm}|L{2.1cm}|L{3.9cm}|L{4.6cm}|C{2.1cm}|}
\caption{Danh mục yêu cầu chức năng và cơ chế kiểm chứng trong hệ thống}\label{tab:g-functional}\\
\hline
\thead{Mã} & \thead{Phân hệ} & \thead{Mô tả yêu cầu nghiệp vụ} & \thead{Cơ chế hiện thực \& Kiểm chứng} & \thead{Trạng thái}\\
\hline
\endfirsthead

\hline
\thead{Mã} & \thead{Phân hệ} & \thead{Mô tả yêu cầu nghiệp vụ} & \thead{Cơ chế hiện thực \& Kiểm chứng} & \thead{Trạng thái}\\
\hline
\endhead

\hline
\endfoot

\hline
\endlastfoot

FR01 & Xác thực & Đăng nhập bằng tài khoản/mật khẩu; băm Argon2id; ký phiên HMAC và hủy phiên tức thì khi đổi mật khẩu/khóa user. & Kiểm tra \texttt{is\_active = 1} và đối chiếu \texttt{auth\_version}; bảo vệ mọi form/API bằng CSRF token 2 lớp. & ĐÃ KIỂM THỬ\\
\hline
FR02 & Phân quyền & Kiểm soát truy cập theo 4 vai trò (RBAC) từ một nguồn duy nhất và cô lập phiếu theo Kỹ thuật viên. & Định nghĩa tập trung tại \texttt{app/services/\allowbreak security\_matrix.py}; chặn vượt quyền dọc/ngang với HTTP 403. & ĐÃ KIỂM THỬ\\
\hline
FR03 & Khách hàng & Quản lý hồ sơ khách hàng (\texttt{customers}): thêm, sửa, xóa có kiểm tra khóa ngoại, tìm kiếm \& phân trang. & Ràng buộc \texttt{phone UNIQUE}; \texttt{ON DELETE RESTRICT} chặn xóa khách đã có xe hoặc lịch hẹn. & ĐÃ KIỂM THỬ\\
\hline
FR04 & Phương tiện & Quản lý hồ sơ xe (\texttt{vehicles}) theo biển số, hãng, dòng, năm sản xuất, số khung VIN và số công-tơ-mét (\texttt{odometer}). & Ràng buộc \texttt{plate UNIQUE}, \texttt{vin UNIQUE}, \texttt{year BETWEEN 1900 AND 2100}, \texttt{odometer >= 0}. & ĐÃ KIỂM THỬ\\
\hline
FR05 & Lịch hẹn & Đặt lịch hẹn dịch vụ (\texttt{appointments}) và quản lý 4 trạng thái (\texttt{SCHEDULED}, \texttt{CONFIRMED}, \texttt{CANCELLED}, \texttt{DONE}). & Kiểm tra \texttt{customer\_id} khớp với chủ sở hữu của \texttt{vehicle\_id}; chặn đặt trùng giờ cho cùng một xe. & ĐÃ KIỂM THỬ\\
\hline
FR06 & Tiếp nhận xe & Ghi biên bản tiếp nhận xe (\texttt{intake\_records}) gắn 1--1 với phiếu sửa chữa. & Lưu \texttt{odometer\_at\_intake}, \texttt{fuel\_level}, tình trạng ngoại thất (\texttt{exterior\_condition}) và phụ kiện kèm theo. & ĐÃ KIỂM THỬ\\
\hline
FR07 & Phiếu sửa chữa & Quản lý vòng đời Phiếu sửa chữa (\texttt{repair\_orders}) qua 6 trạng thái, nhật ký kỹ thuật và nghiệm thu. & Chặn mở 2 phiếu đang chạy cho cùng 1 xe trong khối \texttt{BEGIN IMMEDIATE}; chặn hoàn thành khi phiếu chưa có hạng mục. & ĐÃ KIỂM THỬ\\
\hline
FR08 & Báo giá & Lập báo giá chính thức (\texttt{quotations}, \texttt{quotation\_items}) qua 5 trạng thái và áp dụng vào phiếu sửa chữa. & Áp dụng \textit{idempotent} đúng 1 lần qua cột \texttt{applied\_at}/\texttt{applied\_by}; tự động trừ kho qua sổ cái khi áp dụng. & ĐÃ KIỂM THỬ\\
\hline
FR09 & Kho \& Sổ cái & Quản lý danh mục dịch vụ (\texttt{services}), phụ tùng (\texttt{parts}) và sổ cái tồn kho (\texttt{stock\_transactions}). & Bất biến $\texttt{parts.stock} = \sum \texttt{qty}$; tự động hoàn kho (\texttt{RETURN}) khi gỡ hạng mục hoặc hủy phiếu. & ĐÃ KIỂM THỬ\\
\hline
FR10 & Hóa đơn \& Thu & Lập hóa đơn (\texttt{invoices}) cho phiếu đã hoàn thành, thu tiền nhiều đợt (\texttt{payments}) và hủy khoản thu nhầm. & Ràng buộc \texttt{CHECK (paid <= total)}; kiểm tra chống thu vượt nợ trong khối \texttt{BEGIN IMMEDIATE} (HTTP 409). & ĐÃ KIỂM THỬ\\
\hline
FR11 & Bảo hành & Cấp và theo dõi bảo hành theo từng hạng mục dịch vụ hoặc phụ tùng (\texttt{warranties}). & Ràng buộc \texttt{CHECK XOR} giữa \texttt{SERVICE} và \texttt{PART}; chỉ tạo được cho phiếu \texttt{COMPLETED}. & ĐÃ KIỂM THỬ\\
\hline
FR12 & AI01 \& AI02 & AI01 (\texttt{history}) tóm tắt lịch sử xe 4 mục; AI02 (\texttt{explain}) giải thích dịch vụ 3 mục cho khách hàng. & Xử lý an toàn cả ca biên xe mới chưa có lịch sử (\texttt{vehicle\_id=6}) và dịch vụ trống mô tả (\texttt{service\_id=7}). & ĐÃ ĐO KIỂM\\
\hline
FR13 & AI03 Báo giá & AI03 (\texttt{quote}) sinh báo giá nháp JSON từ hạng mục của phiếu sửa chữa hoặc báo giá đang mở. & Lọc chỉ giữ dòng khớp Tên, Số lượng, Đơn giá với CSDL; ép \texttt{total == backend\_total}; tự động fallback \texttt{\_goi\_mock}. & ĐÃ ĐO KIỂM\\
\hline
FR14 & Human Review & Lưu toàn bộ kết quả AI vào bảng \texttt{ai\_results} ở trạng thái chờ duyệt và kiểm duyệt bởi Quản lý. & Khởi tạo \texttt{review\_status='PENDING'}; chỉ \texttt{MANAGER} (\texttt{ai.review}) được duyệt và chặn duyệt lại lần 2. & ĐÃ KIỂM THỬ\\
\hline
\end{longtable}
\end{small}

\section{Danh mục yêu cầu phi chức năng (NFR)}

Hệ thống đáp ứng 6 yêu cầu phi chức năng cốt lõi (NFR01--NFR06):
\begin{itemize}[noitemsep,topsep=2pt]
  \item \textbf{NFR01 (An toàn xác thực, chống CSRF và phân quyền chặt chẽ)}: Không lưu mật khẩu thô; tự động hủy phiên cũ qua \texttt{auth\_version}; kiểm tra CSRF trên mọi thao tác ghi; cô lập Kỹ thuật viên theo phiếu được phân công.
  \item \textbf{NFR02 (Chính xác tuyệt đối về tiền tệ và tồn kho)}: Mọi cột tiền dùng kiểu số nguyên \texttt{INTEGER} (tính bằng đồng VNĐ); \texttt{unit\_price} được chốt cứng tại thời điểm thêm hạng mục để việc đổi giá danh mục không làm sai lệch hóa đơn cũ; đối soát sổ cái kho luôn đạt chênh lệch bằng 0.
  \item \textbf{NFR03 (Hiệu năng và thời gian phản hồi)}: Các thao tác CRUD nội bộ phản hồi dưới 50ms; thời gian gọi mô hình thật trên Groq (\texttt{openai/gpt-oss-120b}) trung bình từ 2,46s đến 3,99s, đáp ứng tốt yêu cầu dưới 10 giây (tiêu chí CK-07).
  \item \textbf{NFR04 (Kiểm soát AI không tự quyết định tiền)}: Kết quả AI chỉ dừng ở văn bản tư vấn trong bảng \texttt{ai\_results}, tuyệt đối không tự chèn hạng mục vào phiếu sửa chữa hay tự tạo hóa đơn.
  \item \textbf{NFR05 (Chịu lỗi và giới hạn tần suất gọi API)}: Bắt toàn bộ ngoại lệ mạng, quá thời gian chờ (\texttt{AI\_TIMEOUT = 15s}), lỗi HTTP 429/500 từ nhà cung cấp AI để tự động chuyển sang \texttt{\_goi\_mock} mà không làm dừng ứng dụng; giới hạn tốc độ bấm nút chạy AI (\texttt{AI\_RATE\_LIMIT = 10} lần / 60 giây).
  \item \textbf{NFR06 (Tự động nâng cấp CSDL và Sao lưu / Khôi phục an toàn)}: Tự động chạy các bản nâng cấp lược đồ \texttt{001} $\rightarrow$ \texttt{005} (\texttt{database/migrations/run\_missing.py}) một cách \textit{idempotent} khi khởi động; hỗ trợ sao lưu trực tuyến và khôi phục kèm kiểm tra \texttt{PRAGMA integrity\_check}.
\end{itemize}

\section{Kiểm soát phân quyền RBAC và Cô lập phiếu theo Kỹ thuật viên (\texttt{duoc\_gan\_cong})}

Hình~\ref{fig:g-repair-seq} minh họa trình tự phối hợp giữa Lễ tân, Kỹ thuật viên, Quản lý, Backend FastAPI, AI Gateway và SQLite từ lúc tiếp nhận xe đến khi nghiệm thu, thanh toán và cấp bảo hành.

\begin{figure}[H]
\centering
\vspace{2pt}
\resizebox{\textwidth}{!}{%
\begin{tikzpicture}[
    font=\sffamily\scriptsize,
    >=Latex,
    part/.style={rectangle, draw=diagblueline, line width=0.9pt, fill=diagbluefill, rounded corners=4pt, text width=2.7cm, minimum height=0.95cm, align=center, font=\sffamily\scriptsize\bfseries, inner sep=4pt},
    msg/.style={-Latex, line width=0.8pt, navy},
    ret/.style={-Latex, dashed, line width=0.8pt, diaggreenline}
]

\node[part] (p1) at (0,0) {Lễ tân / KTV\\(\texttt{le\_tan} / \texttt{ky\_thuat\_vien})};
\node[part, fill=diagamberfill, draw=diagamberline] (p2) at (3.8,0) {Quản lý Xưởng\\(\texttt{quan\_ly})};
\node[part, fill=diagpurplefill, draw=diagpurpleline] (p3) at (7.8,0) {FastAPI Backend\\(\texttt{app/services/})};
\node[part, fill=diaggreenfill, draw=diaggreenline] (p4) at (12.0,0) {AI Gateway\\(\texttt{ai\_gateway.py})};
\node[part, fill=diagslatefill, draw=diagslateline] (p5) at (16.0,0) {SQLite WAL\\(\texttt{garage.db})};

\foreach \x in {0, 3.8, 7.8, 12.0, 16.0} {
    \draw[dashed, diagslateline!70] (\x,-0.5) -- (\x,-8.2);
}

\draw[msg] (0,-1.0) -- node[above, fill=white, inner sep=1.5pt]{1. Mở Phiếu sửa chữa (\texttt{RECEIVED}) + Biên bản tiếp nhận} (7.8,-1.0);
\draw[msg] (7.8,-1.5) -- node[above, fill=white, inner sep=1.5pt]{2. \texttt{INSERT repair\_orders, intake\_records}} (16.0,-1.5);

\draw[msg] (0,-2.2) -- node[above, fill=white, inner sep=1.5pt]{3. Kiểm tra (\texttt{INSPECTING}) \& Báo giá (\texttt{ESTIMATING})} (7.8,-2.2);
\draw[msg] (0,-2.8) -- node[above, fill=white, inner sep=1.5pt]{4. Gọi AI03 \texttt{POST /api/ai/quote} (Báo giá nháp)} (7.8,-2.8);
\draw[msg] (7.8,-3.3) -- node[above, fill=white, inner sep=1.5pt]{5. Đọc CSDL + Gọi LLM + Lọc giá} (12.0,-3.3);
\draw[ret] (12.0,-3.9) -- node[above, fill=white, inner sep=1.5pt]{6. Lưu \texttt{ai\_results (PENDING)}} (16.0,-3.9);

\draw[msg] (3.8,-4.6) -- node[above, fill=white, inner sep=1.5pt]{7. Quản lý duyệt AI (\texttt{APPROVED})} (7.8,-4.6);
\draw[msg] (3.8,-5.2) -- node[above, fill=white, inner sep=1.5pt]{8. Duyệt \& Áp dụng Báo giá 1 lần} (7.8,-5.2);
\draw[msg] (7.8,-5.8) -- node[above, fill=white, inner sep=1.5pt]{9. \texttt{BEGIN IMMEDIATE}: Trừ kho + Ghi \texttt{stock\_transactions (USAGE)}} (16.0,-5.8);

\draw[msg] (0,-6.5) -- node[above, fill=white, inner sep=1.5pt]{10. KTV được gán (\texttt{duoc\_gan\_cong}) sửa chữa (\texttt{REPAIRING})} (7.8,-6.5);
\draw[msg] (3.8,-7.1) -- node[above, fill=white, inner sep=1.5pt]{11. Nghiệm thu hoàn thành (\texttt{COMPLETED})} (7.8,-7.1);
\draw[msg] (7.8,-7.7) -- node[above, fill=white, inner sep=1.5pt]{12. Xuất \texttt{invoices}, thu \texttt{payments (paid <= total)} \& cấp \texttt{warranties}} (16.0,-7.7);

\end{tikzpicture}%
}

\vspace{2pt}
\caption{Biểu đồ tuần tự (Sequence Diagram) quy trình sửa chữa khép kín và duyệt kết quả AI}
\label{fig:g-repair-seq}
\end{figure}

Toàn bộ ma trận phân quyền của hệ thống được định nghĩa tập trung tại một tệp duy nhất là \texttt{app/services/security\_matrix.py} (biến \texttt{PERMISSIONS}) để dùng chung cho cả Backend lẫn giao diện. Đối với các thao tác sửa đổi trên Phiếu sửa chữa, hàm \texttt{\_kiem\_tra\_quyen\_thao\_tac} trong \texttt{app/services/orders.py} kết hợp kiểm tra vai trò và phạm vi phân công:
{\small
\begin{equation}
\begin{aligned}
\text{DuocThaoTac}(\text{user}, \text{phieu}) \iff\ & (\text{user.role} \in \{\texttt{'MANAGER'}, \texttt{'RECEPTIONIST'}\}) \\
&\lor \left(\text{user.role} = \texttt{'TECHNICIAN'} \land \text{phieu.technician\_id} = \text{user.id}\right)
\end{aligned}
\end{equation}
}

Nhờ chốt chặn này, một Kỹ thuật viên chỉ có thể cập nhật trạng thái, thêm/gỡ hạng mục, ghi nhật ký kỹ thuật hoặc chạy AI03 trên đúng những phiếu sửa chữa được giao cho mình; mọi nỗ lực can thiệp vào phiếu của Kỹ thuật viên khác hoặc phiếu chưa phân công đều bị từ chối ngay tại tầng dịch vụ.

\section{Máy trạng thái Phiếu sửa chữa (6 trạng thái) và Quy trình Human Review}

Vòng đời của một Phiếu sửa chữa (\texttt{repair\_orders.status}) và trạng thái duyệt kết quả AI (\texttt{ai\_results.review\_status}) được quy định chặt chẽ tại biến \texttt{CHUYEN\_TRANG\_THAI} trong \texttt{app/services/security\_matrix.py} và thể hiện trên Hình~\ref{fig:g-ro-state}.

\begin{figure}[H]
\centering
\vspace{2pt}
\resizebox{0.98\textwidth}{!}{%
\begin{tikzpicture}[
  font=\sffamily,
  st/.style={rectangle, rounded corners=3pt, draw=black, line width=0.65pt, fill=white, text width=2.8cm, minimum height=0.85cm, align=center, font=\sffamily\scriptsize\bfseries, inner sep=3pt},
  sthead/.style={rectangle, rounded corners=3pt, draw=black, line width=0.75pt, fill=black!12, text width=2.8cm, minimum height=0.85cm, align=center, font=\sffamily\scriptsize\bfseries, inner sep=3pt},
  aist/.style={rectangle, rounded corners=3pt, draw=black, line width=0.7pt, fill=black!6, text width=4.8cm, minimum height=0.85cm, align=center, font=\sffamily\scriptsize\bfseries, inner sep=3pt},
  arr/.style={-{Latex[length=2mm,width=1.4mm]}, line width=0.6pt, draw=black},
  lbl/.style={font=\sffamily\scriptsize, fill=white, inner sep=1.5pt, align=center}
]
  \node[sthead] (s1) at (0, 2.2) {1. DRAFT / INTAKE\\{\normalfont\scriptsize Mới tiếp nhận}};
  \node[st] (s2) at (4.5, 2.2) {2. DIAGNOSING\\{\normalfont\scriptsize Đang chẩn đoán}};
  \node[st] (s3) at (9.0, 2.2) {3. QUOTED\\{\normalfont\scriptsize Đã lập báo giá}};
  \node[sthead] (s4) at (13.5, 2.2) {4. APPROVED\\{\normalfont\scriptsize Khách đồng ý}};

  \node[st] (s5) at (13.5, 0.0) {5. WAITING\_PARTS\\{\normalfont\scriptsize Chờ đặt giữ kho}};
  \node[st] (s6) at (9.0, 0.0) {6. IN\_PROGRESS\\{\normalfont\scriptsize Đang sửa chữa}};
  \node[st] (s7) at (4.5, 0.0) {7. QC\_PASSED\\{\normalfont\scriptsize Nghiệm thu kỹ thuật}};
  \node[sthead] (s8) at (0, 0.0) {8. INVOICED / CLOSED\\{\normalfont\scriptsize Thanh toán \& Đóng}};

  \node[aist] (ai_pend) at (3.2, -1.9) {AI\_SUGGESTION: PENDING\\{\normalfont\scriptsize Chờ cố vấn thẩm định}};
  \node[aist] (ai_acc) at (10.8, -1.9) {AI\_SUGGESTION: ACCEPTED / REJECTED\\{\normalfont\scriptsize Ghi Audit \& Chuyển vào báo giá RO}};

  \draw[arr] (s1) -- node[lbl, above]{Giao KTV} (s2);
  \draw[arr] (s2) -- node[lbl, above]{Lên báo giá} (s3);
  \draw[arr] (s3) -- node[lbl, above]{Duyệt giá} (s4);
  \draw[arr] (s4) -- node[lbl, right]{Kiểm kho} (s5);
  \draw[arr] (s5) -- node[lbl, above]{Xuất kho} (s6);
  \draw[arr] (s6) -- node[lbl, above]{Hoàn tất} (s7);
  \draw[arr] (s7) -- node[lbl, above]{Thu phí} (s8);

  \draw[arr, dashed] (s2.north) -- ++(0, 0.45) -| (-1.95, -1.9) -- node[lbl, above=2pt]{Gọi AI01/AI02} (ai_pend.west);
  \draw[arr] (ai_pend.east) -- node[lbl, above]{Duyệt HITL} (ai_acc.west);
  \draw[arr, dashed] (ai_acc.east) -- (15.4, -1.9) |- (s3.north);
\end{tikzpicture}%
}

\vspace{2pt}
\caption{Máy trạng thái vòng đời Phiếu sửa chữa (6 trạng thái) và Trạng thái duyệt kết quả AI}
\label{fig:g-ro-state}
\end{figure}

Bốn quy tắc bảo vệ trạng thái quan trọng được cài đặt trong \texttt{app/services/orders.py} và \texttt{app/services/ai.py} gồm:
\begin{enumerate}[noitemsep,topsep=2pt]
  \item \textbf{Lộ trình chuyển trạng thái chống nhảy cóc (\texttt{CHUYEN\_TRANG\_THAI})}:
  \begin{itemize}[noitemsep,topsep=1pt]
    \item Từ \texttt{RECEIVED} chỉ được chuyển sang \texttt{INSPECTING} hoặc \texttt{CANCELLED}.
    \item Từ \texttt{INSPECTING} chỉ được chuyển sang \texttt{ESTIMATING} hoặc \texttt{CANCELLED}.
    \item Từ \texttt{ESTIMATING} được chuyển sang \texttt{REPAIRING}, quay lại \texttt{INSPECTING} hoặc hủy (\texttt{CANCELLED}).
    \item Từ \texttt{REPAIRING} được chuyển sang \texttt{COMPLETED} (khi đã có ít nhất 1 hạng mục) hoặc quay lại \texttt{INSPECTING} khi phát sinh hạng mục cần kiểm tra lại.
    \item Hai trạng thái \texttt{COMPLETED} và \texttt{CANCELLED} là trạng thái đóng (\texttt{[]}), đồng thời khóa hoàn toàn việc thêm/gỡ hạng mục (\texttt{TRANG\_THAI\_KHOA\_HANG\_MUC}).
  \end{itemize}
  \item \textbf{Khóa hạng mục khi đã phát hành hóa đơn}: Hàm \texttt{\_khoa\_hang\_muc} kiểm tra nếu phiếu đã có bản ghi trong bảng \texttt{invoices} thì chặn mọi thao tác thêm hoặc gỡ hạng mục để tổng tiền hạng mục và tổng tiền hóa đơn không bao giờ lệch nhau.
  \item \textbf{Tự động hoàn trả phụ tùng về kho khi gỡ hạng mục hoặc hủy phiếu}: Khi gỡ một dòng phụ tùng hoặc chuyển phiếu sang \texttt{CANCELLED}, hệ thống tự động gọi \texttt{ghi\_giao\_dich(..., +qty, 'RETURN')} trong cùng một giao dịch \texttt{BEGIN IMMEDIATE} để cộng trả lại \texttt{parts.stock} và ghi dòng đối ứng vào \texttt{stock\_transactions}.
  \item \textbf{Duyệt kết quả AI đúng một lần (Human Review)}: Kết quả AI sinh ra luôn ở trạng thái \texttt{PENDING}; chỉ tài khoản có quyền \texttt{ai.review} (\texttt{MANAGER}) mới được chuyển sang \texttt{APPROVED} hoặc \texttt{REJECTED}, và khi đã duyệt thì hệ thống chặn duyệt lại lần thứ hai.
\end{enumerate}

\section{Kết luận chương}

Chương 3 đã đặc tả đầy đủ 4 vai trò nhân sự, 14 yêu cầu chức năng, 6 yêu cầu phi chức năng, cơ chế cô lập phiếu của Kỹ thuật viên (\texttt{duoc\_gan\_cong}) và máy trạng thái 6 bước của Phiếu sửa chữa. Đây là cơ sở trực tiếp cho thiết kế kiến trúc và cơ sở dữ liệu ở Chương 4.
